%!TEX TS-program = lualatex
%!TEX encoding = UTF-8 Unicode

\documentclass[letterpaper]{tufte-handout}

%\geometry{showframe} % display margins for debugging page layout

\usepackage{fontspec}
\def\mainfont{Linux Libertine O}
\setmainfont[Ligatures={Common,TeX}, Contextuals={NoAlternate}, BoldFont={* Bold}, ItalicFont={* Italic}, Numbers={OldStyle}]{\mainfont}
\setsansfont[Scale=MatchLowercase, Numbers={OldStyle}]{Linux Biolinum O} 
\setmonofont{Linux Libertine O}
\usepackage{microtype}

\usepackage{graphicx} % allow embedded images
  \setkeys{Gin}{width=\linewidth,totalheight=\textheight,keepaspectratio}
  \graphicspath{{/Users/goby/Documents/teach/434/lectures/}} % set of paths to search for images

\usepackage{amsmath}  % extended mathematics
\usepackage{booktabs} % book-quality tables
\usepackage{units}    % non-stacked fractions and better unit spacing
\usepackage{multicol} % multiple column layout facilities
%\usepackage{fancyvrb} % extended verbatim environments
%  \fvset{fontsize=\normalsize}% default font size for fancy-verbatim environments

\usepackage{enumitem}
\usepackage{mhchem}

\makeatletter
% Paragraph indentation and separation for normal text
\renewcommand{\@tufte@reset@par}{%
  \setlength{\RaggedRightParindent}{1.0pc}%
  \setlength{\JustifyingParindent}{1.0pc}%
  \setlength{\parindent}{1pc}%
  \setlength{\parskip}{0pt}%
}
\@tufte@reset@par

% Paragraph indentation and separation for marginal text
\renewcommand{\@tufte@margin@par}{%
  \setlength{\RaggedRightParindent}{0pt}%
  \setlength{\JustifyingParindent}{0.5pc}%
  \setlength{\parindent}{0.5pc}%
  \setlength{\parskip}{0pt}%
}
\makeatother

% Set up the spacing using fontspec features
   \renewcommand\allcapsspacing[1]{{\addfontfeatures{LetterSpace=15}#1}}
   \renewcommand\smallcapsspacing[1]{{\addfontfeatures{LetterSpace=10}#1}}

\title{{\scshape zo} 478/678 Study Guide 02}

\date{} % without \date command, current date is supplied

%\newcommand\lecturefileA{478_lecture01_instructor}
%\newcommand\lecturefileB{478_lecture02_instructor}

\begin{document}

\maketitle	% this prints the handout title, author, and date

%\printclassoptions

\section{Vocabulary}\marginnote{\textbf{Study:} pgs. 113--119, 57--64. Webb, P. 1984. Form and function in fish swimming. Scientific American 251: 72--82. Emphasize the mechanisms of swimming but also review form and function from Lecture 01.} 
\vspace{-1\baselineskip}
\begin{multicols}{2}
myomere\\
hypaxial\\
epaxial\\
undulation \\
oscillation \\
anguilliform swimming \\
subcarangiform swimming \\
carangiform swimming \\
ostraciform swimming \\
rajiform swimming \\
amiiform swimming  \\
labriform swimming \\
balistiform swimming \\
gymnotiform swimming \\
tetraodontform swimming \\
normoxic \\
hypoxic \\
anoxic \\
gill rakers \\
gill arch \\
gill filaments \\
lamellae (singular: lamella) \\
countercurrent exchange \\
buccal cavity \\
opercular cavity \\
ram ventilation \\
suprabranchial organ \\
labyrinth organ \\
sinus venosus \\
atrium \\
ventricle \\
conus arteriosus \\
bulbus arteriosus \\
ventral aorta \\
dorsal aorta \\
afferent branchial arteries \\
efferent branchial arteries \\
common cardinal vein \\
hemoglobin \\
tetrameric \\ \hspace{2ex}hemoglobin \\
oxygen dissociation curve \\
P$_{50}$ \\
Bohr effect$^1$ \\
Root effect$^1$
\end{multicols}
\marginnote[-3\baselineskip]{$^1$ Remember ABCR: Affinity\,=\,Bohr, Capacity\,=\,Root. Or, you can BARC like a dogfish: Bohr\,=\,Affinity, Root\,=\,Capacity.}

\section{Concepts}

\begin{enumerate}
	\item Briefly explain the function of myomeres for swimming in fishes. Draw an illustration of a myomere in relation to the body of the fish. 
	
	\item Explain how myomere contraction and relaxation results in undulation of the body.
	
	\item Be able to differentiate each of the types of swimming listed above.  Which use undulation?  Which use oscillation? Relate this to the assigned reading (Webb, Form and Function) that I handed out in class.
	
	\item List the swimming types that used the median or paired fins. For each, state whether it is based on undulation or oscillation.
	
	\item List the swimming types that uses body or caudal fins. For each, state whether it is based on undulation or oscillation.
	
	\item How much dissolved oxygen (\%) can pure, cold freshwater hold?  What other factors influence how much oxygen can be dissolved in water.
	
	\item Describe the specific mechanisms in the gill that fishes use to extract as much oxygen as possible from the water.
	
	\item Describe the relationship between gill surface area and fish form and function. What mechanisms are used to increase the surface area of the gills?
	
	\item Describe alternative physical and physiological mechanisms that some fishes utilize to obtain oxygen in normal and low oxygen conditions.
	
	\item Describe and illustrate countercurrent exchange in the gills.  Explain how countercurrent exchange maximizes oxygen uptake.
	
	\item Describe the mechanics of gill ventilation (pumping water across the gills).
	
	\item \textbf{Thoughtful:} Despite the fact that O$_2$ has slightly lower solubility in saltwater than in freshwater, most of the air-breathing fishes are freshwater species.  Why do you think this is? 

	\item Which chamber of the heart provides most acceleration of the blood? 

	\item What are the functions of the sinus venosus, and the conus / bulbus arteriosus?
	
	\item Describe and illustrate how the circulatory/respiratory system differs between lungfishes and teleost fishes.
	
	\item How is the heart of the lungfish structured to keep oxygenated and deoxygenated blood separate from each other?
	
	\item Interpret and draw blood oxygen equilibrium curves (oxygen affinity curves).  Compare oxygen affinity between an active fish in normoxic water vs a sluggish fish in hypoxic water.
	
	\item Describe the effects that pH and concentration of dissolved CO$_2$ have on the oxygen affinity of hemoglobin.
	
	\item Explain how the Bohr Effect and the oxygen affinity of hemoglobin influence the loading and unloading of oxygen at the gills and in body tissue.
	
	\item Be able to write and interpret the bicarbonate buffering equilibrium equation. Explain how this system (the equilibrium equation) contributes to the loading and unloading of oxygen at the gills and in body tissue.  
\end{enumerate}


\end{document}